% ---- pg-288 （书页 276） ----

\begin{center}
\small 续表
\end{center}

\begin{mvtable}
\begin{tabular}{rrccr}
\toprule
560 纳米 & 720 纳米 & 品种 & 时间 & 复制\\\midrule
8.73 & 23.27 & LP & 1 & 1\\
7.94 & 20.87 & LP & 1 & 2\\
8.37 & 22.16 & LP & 1 & 3\\
7.86 & 21.78 & LP & 1 & 4\\
8.45 & 26.32 & LP & 2 & 1\\
6.79 & 22.73 & LP & 2 & 2\\
8.34 & 26.67 & LP & 2 & 3\\
7.54 & 24.87 & LP & 2 & 4\\
14.04 & 44.44 & LP & 3 & 1\\
13.51 & 37.93 & LP & 3 & 2\\
13.33 & 37.93 & LP & 3 & 3\\
12.77 & 60.87 & LP & 3 & 4\\
\bottomrule
\end{tabular}
\end{mvtable}

\vspace{2pt}
\noindent\rule{\linewidth}{0.4pt}

\noindent\small 资料来源: 数据经 Mairtin Mac Siurtain 许可使用.

\vspace{1em}
\noindent 6.34\quad 参看例 6.15.

\noindent (a) 在同一图上分别画出 $\bar{x}_1$ 的分量和 $\bar{x}_2$ 的分量对时间的轮廓, 进行比较, 并对比较结果作出评论.

\noindent (b) 取 $\alpha=0.01$, 检验线性生长假设是否合适?

\noindent 6.35\quad 参看例 6.15, 但要把所有的 31 个对象当作一个组来处理. $(q+1)\times 1$ 阶向量 $\boldsymbol{\beta}$ 的极大似然估计量为: $\hat{\boldsymbol{\beta}}=(\mathbf{B}'\mathbf{S}^{-1}\mathbf{B})^{-1}\mathbf{B}'\mathbf{S}^{-1}\bar{\mathbf{x}}$, 其中 $\mathbf{S}$ 为样本协方差阵.

\noindent 极大似然估计的协方差的估计量为
\[
\widehat{\operatorname{Cov}}(\hat{\boldsymbol{\beta}})=\frac{(n-1)(n-2)}{(n-1-p+q)(n-p+q)n}(\mathbf{B}'\mathbf{S}^{-1}\mathbf{B})^{-1}
\]
用二次生长曲线拟合这组数据, 并加以讨论.

\noindent 6.36\quad 参考例 6.4. 在这个例子中, 给出用电量的汇总信息, 采用博克斯 $M$ 检验来检验零假设 $H_0: \boldsymbol{\Sigma}_1=\boldsymbol{\Sigma}_2=\boldsymbol{\Sigma}$, 这里 $\boldsymbol{\Sigma}_1$ 是威斯康星州有空调用户总体的消耗量的两个观测值的协方差矩阵, $\boldsymbol{\Sigma}_2$ 是威斯康星州没有空调用户总体的消耗量的两个观测值的协方差矩阵. 取 $\alpha=0.05$.

\noindent 6.37\quad \hyperref[tab:6.9]{表 6.9} 中包含了 24 个雌性和 24 个雄性龟的甲壳测量数据. 采用博克斯 $M$ 检验来检验零假设 $H_0: \boldsymbol{\Sigma}_1=\boldsymbol{\Sigma}_2=\boldsymbol{\Sigma}$. 这里 $\boldsymbol{\Sigma}_1$ 是雌性龟甲壳测量数据的总体协方差矩阵, $\boldsymbol{\Sigma}_2$ 是雄性龟甲壳测量数据的总体协方差矩阵. 取 $\alpha=0.05$.

\noindent 6.38\quad \hyperref[tab:11.7]{表 11.7} 中包含了来自三个砂岩地区的原油样本中的三种痕量元素和两种烃类化合物含量的观测值. 采用博克斯 $M$ 检验来检验三个地区的总体协方差矩阵是否相等. 取 $\alpha=0.05$. 这里有 $p=5$ 个变量, 你可以考虑对这些变量的观测值进行变换使得它们更接近正态分布.

\noindent 6.39\quad 水蟒是世界上最大的蛇类之一. 研究者捕捉水蟒并测量 (i) 嘴到排泄口的长度; (ii) 重量. 这些测量数据见\hyperref[tab:6.19]{表 6.19}.

\noindent (a) 取 $\alpha=0.05$, 应用大样本统计检验雌性和雄性水蟒总体的均值是否相等.

\noindent (b) 在这种情形下, 融合方差是否合理? 请给出原因.

\noindent (c) 对雌性和雄性水蟒的长度和体重的均值差, 构造其 95\% 庞弗罗尼置信区间.
